% problem-000321 4 拟卤素氰化学
\section*{4. 拟卤素氰化学}
\textit{下列为分问。}

\subsection*{4-1}
$\mathrm { H g C l _ { 2 } }$ 和 $\mathrm { H g ( C N ) _ { 2 } }$ 反应可制得(CN)_{2}，写出反应⽅程式。

\subsection*{4-2}
画出 $\mathrm { C N ^ { - } }$ $( \mathrm { C N } ) _ { 2 }$ 的路易斯结构式。

\subsection*{4-3}
写出 $\mathrm { ( C N ) _ { 2 } ( g ) }$ 在 $\mathrm { O _ { 2 } ( g ) }$ 中燃烧的反应⽅程式。

\subsection*{4-4}
298 K 下，(CN)_{2}(g)的标准摩尔燃烧热为−1095 kJ $\mathrm { m o l ^ { - 1 } }$ $\mathrm { C _ { 2 } H _ { 2 } ( g ) }$ 的标准摩尔燃烧热为−1300 kJ $\mathrm { m o l ^ { - 1 } }$ $\mathrm { C _ { 2 } H _ { 2 } ( g ) }$ 的标准摩尔⽣成焓为227 kJ $\mathrm { m o l ^ { - 1 } }$ $\mathrm { H _ { 2 } O ( l ) }$ 的标准摩尔⽣成焓为-286 kJ $\mathrm { m o l ^ { - 1 } }$ ，计算 $\mathrm { ( C N ) _ { 2 } ( g ) }$ 的标准摩尔⽣成焓。

\subsection*{4-5}
(CN)_{2} 在 $3 0 0 { \sim } 5 0 0 \ ^ { \circ } \mathrm { C }$ 形成具有⼀维双链结构的聚合物，画出该聚合物的结构。

\subsection*{4-6}
电镀⼚向含氰化物的电镀废液中加⼊漂⽩粉以消除有毒的CN-，写出化学⽅程式（漂⽩粉⽤ClO-表⽰）。
